\chapter{TWA 的有效范围与失效模式}
\label{ch:validity}

\begin{learningobjectives}
完成本章后，读者应能：
\begin{enumerate}
  \item 把 TWA 误差分解为状态表示、动力学截断、数值离散和抽样误差；
  \item 根据占据、模式数、不稳定性、观测量和时间窗口评估风险；
  \item 设计模式截止、轨迹数和独立方法的验证矩阵；
  \item 避免用经典场晚时热化替代量子热化结论。
\end{enumerate}
\end{learningobjectives}

\section{有效性不是一个布尔变量}

“TWA 对这个模型有效吗”缺少三个必要限定：对哪个初态、哪个观测量、哪个时间窗口。第5章的二次系统对任意时间和任意初态分布都有精确 Wigner 漂移；第6章 Kerr 模型在早期相干性上准确，却完全错过回复；第8章二聚体在指定参数和数差观测量上通过有限时间基准。这些例子说明，有效性必须写成带条件的命题。

一项 TWA 结果至少包含四层近似：
\begin{enumerate}
  \item \textbf{状态表示}：初态 Wigner 函数是否精确可抽样，或只匹配了有限矩；
  \item \textbf{动力学截断}：高阶 Moyal 导数是否在目标窗口内可忽略；
  \item \textbf{数值表示}：有限模式、网格、时间步长和浮点误差；
  \item \textbf{统计估计}：有限轨迹和相关样本导致的不确定度。
\end{enumerate}
只有第四层误差必然按 $N_s^{-1/2}$ 缩小。增加轨迹数对前三层没有系统修复作用。

\section{支持 TWA 的条件}

\subsection{相关模式高占据}

半经典展开通常要求对动力学和观测量重要的模式具有大于单位量级的占据。这里的“相关模式”不是任意选定的粗网格模式；若非线性散射不断把半量子噪声送入低占据高能模，总粒子数很大仍不足以保证可靠。经典场研究常用“粒子数显著大于保留模式数”作为必要的经验条件之一\cite{Sinatra2002,Blakie2008}。

\subsection{有限时间和光滑相空间结构}

初始 Wigner 分布在相空间尺度上光滑，且经典流尚未形成细密折叠时，高阶导数较不重要。混沌或动力学不稳定可以指数放大微小结构，使有效时间明显短于平稳系统。不存在只依赖 $N$、适用于所有模型和观测量的统一失效时间公式。

\subsection{低阶、粗粒化观测量}

密度、低阶等时关联和大尺度结构因子通常更适合 TWA。Wigner 负值、高阶尾部、量子回波和精确回复对干涉结构更敏感；即使平均占据很高，相关计算也可能失败。

\section{模式截止与半量子噪声}

每个保留玻色模式在真空中贡献半个对称排序量子。若保留 $M$ 个模式，初始经典场中含有数量级为 $M/2$ 的虚拟占据。正确的等时估计器会减去对应排序项，但非线性轨迹仍会让这些高能噪声参与散射。

因此，细化网格会同时做两件事：提高空间分辨率，并加入更多真空噪声模式。若结果随截止显著变化，不能简单选择“最细网格”作为答案。需要寻找物理观测量对截止稳定的窗口，或采用正规化、投影经典场和更合适的量子方法。针对凝聚气体的分析表明，晚时经典场可能趋向与目标量子温度不同的经典平衡，从而产生伪阻尼\cite{Sinatra2002}。

\begin{warningbox}
网格收敛不是单调的“越细越好”。在 TWA 中，改变网格还改变了真空噪声总量和可发生散射的高能通道；必须同时记录模式数、粒子数和截止能量。
\end{warningbox}

\section{长时间热化的特别风险}

闭合量子系统的幺正演化微观可逆。局域观测量可以因去相位和模式混合趋于稳态，但“不可逆”通常是粗粒化意义。TWA 轨迹同样是可逆 Hamilton 流，其系综可以相空间混合；与此同时，经典非线性场可能向 Rayleigh--Jeans 型分布演化。后者依赖紫外截止，不能自动等同于 Bose--Einstein 量子热平衡。

声称“完全热化”至少应提供：
\begin{enumerate}
  \item 能量、粒子数和其他守恒量的闭合账目；
  \item 多个初态在相同守恒量下趋于同一粗粒化结果；
  \item 模式分布与候选统计系综的定量比较；
  \item 模式截止和网格变化下的稳定性；
  \item 在可重叠区与动理学、精确小系统或其他量子方法比较；
  \item 对预热化、积分性和有限尺寸回复的排除或限定。
\end{enumerate}
仅观察相干因子下降或单个动量峰变宽，不足以满足这些条件。

\section{验证矩阵}

推荐把验证设计成互补矩阵：
\begin{center}
\begin{tabularx}{\textwidth}{p{0.20\textwidth}X X}
\toprule
检查 & 能证明什么 & 不能证明什么 \\
\midrule
增加轨迹数 & Monte Carlo 均值和误差条稳定 & 动力学截断正确 \\
减小时间步长 & 积分器收敛、守恒误差下降 & 连续 TWA 接近量子解 \\
改变网格/截止 & 结果对数值表示稳定 & 所有未保留模式都不重要 \\
二次极限 & 抽样、漂移和排序实现一致 & 非线性长时间可靠 \\
精确小系统 & 指定参数和观测量的总误差 & 热力学极限必然相同 \\
Bogoliubov 极限 & 弱涨落线性响应正确 & 强非线性晚时正确 \\
\bottomrule
\end{tabularx}
\end{center}

最有价值的验证不是重复同一类收敛，而是让不同误差源具有不同响应。例如轨迹数变化只影响统计噪声，时间步长变化影响守恒漂移，Fock 截断影响精确基准，系统尺寸变化影响有限尺寸物理。

\section{按风险选择替代方法}

若初态具有强 Wigner 负值，可考虑 positive-P、离散 Wigner 或直接 Fock/张量网络表示；若系统是一维低纠缠短程模型，MPS 往往比连续场 TWA 更受控；若占据低但模式数少，精确对角化更直接；若关注弱耦合长时间碰撞，量子 Boltzmann 或 Keldysh 动理学可能比盲目延长 TWA 更合适。开放系统则需保留主方程产生的扩散和噪声。

方法选择不应以“哪个能跑到更大尺寸”为唯一标准。还要问目标结论依赖哪类量子信息，以及候选方法是否保存这些信息。

\section{结论强度与证据强度}

建议采用以下层级：
\begin{enumerate}
  \item \textbf{数值观察}：在给定离散化和轨迹数下出现某趋势；
  \item \textbf{TWA 稳健结果}：趋势通过抽样、步长与截止检查；
  \item \textbf{受基准支持的量子结论}：在重叠区与独立量子方法一致；
  \item \textbf{热力学或普适结论}：还需系统尺寸标度、序参量定义和理论论证。
\end{enumerate}
写作时应把结论停在证据实际达到的层级。第17章讨论超固态时将严格使用这一规则。

\begin{chaptercheck}
在运行大规模轨迹前先写下失效判据：哪些模式可能变成低占据，哪个观测量最敏感，何时与基准比较，怎样改变截止。若只有“轨迹数足够多”这一项，验证计划尚未完成。
\end{chaptercheck}

\section*{本章小结}

TWA 的有效性依赖初态、相关模式占据、动力学稳定性、观测量和时间窗口。半量子噪声使模式截止成为物理与数值耦合的问题，尤其限制晚时热化解释。可靠结论需要把抽样、积分、表示和方法误差分别检验，并用与结论强度相匹配的独立证据支撑。

\begin{exercises}
  \item 为第6章 Kerr 例子填写本章验证矩阵，并指出哪个检查能暴露回复失败。
  \item 设一维网格长度固定、格点数加倍，计算真空半量子总数怎样变化。
  \item 设计一个区分相位混合与能量热化的数值协议。
  \item 对一个含 $10^5$ 个粒子和 $10^4$ 个模式的模拟，讨论总平均占据为何不足以描述占据分布风险。
  \item 为“系统在热力学极限中形成长程密度序”列出至少四类独立证据。
\end{exercises}
